\section{Tổng quan}
\label{sec:tong-quan}

\subsection{Bài toán}
Đề bài yêu cầu xây dựng một ứng dụng cộng tác theo mô hình client--server hoạt động trong mạng nội bộ (LAN).
Người dùng có thể đăng nhập, xem danh sách những người đang online, chat chung với cả phòng (broadcast), nhắn
tin riêng cho từng người (unicast) và gửi file cho nhau.

Yêu cầu trọng tâm và cũng là phần thử thách nhất của đề là toàn bộ dữ liệu phải truyền qua \textbf{một giao thức
tầng ứng dụng do nhóm tự thiết kế trên nền UDP}. Vì UDP vốn chỉ gửi từng gói tin (datagram) rời rạc, không đảm
bảo gói tới đích và cũng không giữ thứ tự, nhóm cần tự xây dựng thêm các cơ chế ở tầng trên để đáp ứng:
\begin{itemize}
  \item Truyền dữ liệu tin cậy và đúng thứ tự;
  \item Flow control -- tránh gửi nhanh hơn khả năng xử lý của bên nhận và Congestion control -- điều chỉnh tốc độ gửi theo tình trạng đường truyền, cài đặt thuật toán CUBIC \cite{rfc9438} hoặc BBR \cite{bbr};
  \item Mã hóa và xác thực dữ liệu trên đường truyền;
  \item Loại bỏ \emph{Head-of-Line blocking} (HoL), tức là khi đang truyền một file lớn thì các tin nhắn chat gửi xen kẽ không bị kẹt lại chờ file truyền xong.
\end{itemize}

\subsection{Lựa chọn công nghệ}
Nhóm hiện thực toàn bộ hệ thống bằng ngôn ngữ \textbf{Go} và \textbf{chỉ sử dụng thư viện chuẩn} (standard
library). Các package chính được dùng gồm:
\begin{itemize}
  \item \fname{net}: mở và đọc/ghi trực tiếp trên socket TCP, UDP thô (\fname{net.ListenUDP},
        \fname{WriteToUDP}, \fname{ReadFromUDP}).
  \item \fname{crypto/ecdh}, \fname{crypto/aes}, \fname{crypto/cipher}, \fname{crypto/sha256}: thực hiện trao đổi
        khóa X25519 \cite{rfc7748} và mã hóa đối xứng AES-256-GCM \cite{gcm}.
  \item \fname{encoding/binary}, \fname{encoding/json}: đóng gói header/frame nhị phân và serialize các bản tin
        tầng ứng dụng dưới dạng JSON.
\end{itemize}
Lý do chính nhóm chọn Go là vì cơ chế goroutine rất nhẹ (mỗi goroutine lúc khởi tạo chỉ tốn vài KB stack). Nhờ đó nhóm có thể tách riêng các luồng đọc, ghi và quản lý trạng thái cho từng kết nối; ngay cả khi server chạy đồng thời 20 client thì tổng số goroutine cũng chỉ khoảng vài chục. Ngoài ra, cấu trúc \fname{select} kết hợp với channel của Go rất hợp để viết event loop cho giao thức mạng mà hạn chế phải dùng quá nhiều khóa \texttt{Mutex}, giảm hẳn nguy cơ deadlock hay race condition.

\subsection{Phân tầng}
Để đáp ứng yêu cầu tách bạch giữa xử lý mạng, logic nghiệp vụ và giao diện, nhóm chia mã nguồn thành bốn tầng như ở \cref{tab:tang}.

\begin{table}[H]
  \centering
  \caption{Các tầng của hệ thống}
  \label{tab:tang}
  \begin{tabularx}{\linewidth}{lL{0.3\linewidth}Y}
    \toprule
    \textbf{Tầng} & \textbf{Package} & \textbf{Nhiệm vụ} \\
    \midrule
    Giao diện & \fname{internal/webui} & HTTP/1.1 cài đặt trên socket TCP, phục vụ trang HTML + JS thuần \\
    Ứng dụng & \fname{internal/server}, \fname{internal/client} & Xác thực đăng nhập, quản lý phiên online, định
    tuyến chat/DM, chuyển tiếp file \\
    Định dạng & \fname{internal/protocol} & Định nghĩa header gói, các loại frame, bản tin ứng dụng (\fname{AppMsg},
    \fname{FileMeta}) và bộ đệm sắp thứ tự \fname{SeqBuffer} \\
    Transport & \fname{internal/transport} & Truyền tin cậy trên UDP, mã hóa, ghép kênh nhiều stream độc lập,
    điều khiển luồng và tắc nghẽn \\
    \bottomrule
  \end{tabularx}
\end{table}

Giữa tầng ứng dụng và tầng transport, nhóm giao tiếp thông qua một interface chung là \fname{transport.Endpoint}
(\cref{sec:api}). Nhờ lớp trừu tượng này, phần code ở \fname{server} và \fname{client} chỉ cần gọi hàm gửi một
bản tin ngắn, mở một stream dài để truyền file, hoặc đọc các mảnh byte đã được sắp xếp đúng thứ tự từ channel
nhận về. Toàn bộ phần việc nặng ở dưới như chia nhỏ gói tin, đánh số thứ tự, gửi ACK, truyền lại gói mất hay mã
hóa đều được gói gọn bên trong \fname{internal/transport}.

\subsection{Đối chiếu với 10 bước của đề}
\label{subsec:rubric}

\begin{longtable}{c L{0.2\linewidth} L{0.42\linewidth} L{0.2\linewidth}}
  \caption{Đối chiếu 10 bước của đề với thiết kế}\label{tab:rubric} \\
  \toprule
  \textbf{Bước} & \textbf{Yêu cầu} & \textbf{Cách đáp ứng} & \textbf{Ở đâu} \\
  \midrule
  \endfirsthead
  \toprule
  \textbf{Bước} & \textbf{Yêu cầu} & \textbf{Cách đáp ứng} & \textbf{Ở đâu} \\
  \midrule
  \endhead
  1 & Thiết kế giao thức & Gói UDP $\le$ 1200 byte, 4 loại frame, ghép kênh nhiều stream trong một kết nối,
  phân loại stream MESSAGE và FILE & \cref{sec:giao-thuc} \\
  2 & Handshake TCP & Chuỗi bản tin \texttt{HELLO} $\to$ \texttt{KEYX} $\to$ \texttt{AUTH} $\to$
  \texttt{AUTH\_OK}/\texttt{AUTH\_FAIL} đóng khung \texttt{[len][JSON]}; mỗi phiên handshake chạy trên một
  goroutine riêng với deadline 10 s & \cref{subsec:handshake}, \cref{fig:seq-login} \\
  3 & Đăng nhập, quản lý trạng thái & File \texttt{users.json} chỉ lưu \texttt{sha256(salt $\|$ pass)};
  \fname{Registry} dùng \texttt{RWMutex} tra cứu session theo username và \texttt{conn\_id}, chặn đăng nhập
  trùng ngay trong lúc giữ khóa; server tự ghi đè trường \texttt{From} để chống giả mạo; tự động broadcast lại
  danh sách online khi có người vào/ra & \cref{sec:kien-truc}, \cref{sec:bao-mat} \\
  4 & Concurrency & Transport: mỗi kết nối có đúng 1 goroutine giữ state, giao tiếp qua channel. Server: mỗi
  session tách riêng goroutine đọc (\fname{serveSession}) và goroutine ghi (\fname{writeLoop}) kèm hàng đợi
  riêng, tránh việc một client đọc chậm làm nghẽn cả server; chạy ổn định với 20 client đồng thời
  (\cref{tab:he-thong}) & \cref{subsec:concurrency}, \cref{sec:kien-truc} \\
  5 & Streaming nhạy thời gian & Mỗi tin nhắn đi trên 1 stream ưu tiên cao riêng biệt; loại bỏ HoL ở cả chiều
  gửi lẫn chiều nhận & \cref{sec:hol} \\
  6 & Truyền file tin cậy qua UDP & Cơ chế SACK, phát hiện mất gói theo 3 điều kiện, chỉ truyền lại đúng frame bị
  mất, ghép dữ liệu theo offset, kiểm tra toàn vẹn bằng SHA-256 đầu cuối & \cref{sec:tin-cay} \\
  7 & Resilience & Keep-alive 2 s, idle timeout 10 s, frame \texttt{CLOSE}, bỏ qua gói lỗi/giả mạo,
  RTO exponential backoff. Server bọc \texttt{recover()} cho từng session và dọn dẹp session an toàn. Client hỗ
  trợ tự kết nối lại (backoff 1--30 s) và bắt tín hiệu Ctrl+C để gửi \texttt{CLOSE} trước khi thoát
  & \cref{subsec:song-chet}, \cref{subsec:state} \\
  8 & Bảo mật & Trao đổi khóa ECDH X25519 kết hợp AES-256-GCM, gắn header làm AAD, đảm bảo nonce không trùng lặp
  & \cref{sec:bao-mat} \\
  9 & Giao diện HTTP + JS thuần & Tự cài đặt parser HTTP/1.1 trực tiếp trên \fname{net.Listen} (xử lý request
  line, header, \texttt{Content-Length}); giao diện HTML + JS thuần dùng cơ chế polling
  \texttt{/api/events?since=N}; upload file qua \texttt{POST} raw body; nhúng file tĩnh vào binary bằng
  \texttt{embed} & \cref{sec:kien-truc} \\
  10 & Test, chất lượng code & Kiểm tra \fname{go vet} và \fname{go test -race}; viết unit test cho
  protocol/transport, integration test cho server (\fname{FakePair}) và end-to-end test qua UDP thật với loss
  10\%; đo đạc hiệu năng trên bộ giả lập mạng & \cref{sec:kiem-thu} \\
  \bottomrule
\end{longtable}

